\documentclass[11pt, a4paper, twoside]{article}

% --- Packages ---
\usepackage[utf8]{inputenc}
\usepackage{amsmath, amssymb, amsthm}
\usepackage{geometry}
\usepackage{tcolorbox}
\usepackage{fancyhdr}

% --- Page Layout & Styling ---
\geometry{
	top=1in, 
	bottom=1in, 
	left=1.25in, 
	right=1.25in
}

\pagestyle{fancy}
\fancyhf{}
\fancyhead[LE,RO]{\thepage}
\fancyhead[LO]{\small \textit{Practice Problems: Advanced Coding Theory}}
\fancyhead[RE]{\small \textit{Galois Descent \& Classification}}
\renewcommand{\headrulewidth}{0.4pt}

% Custom boxes
\newtcolorbox{problembox}[1][]{
	colback=blue!5!white, colframe=blue!75!black,
	fonttitle=\bfseries, title=#1, arc=3pt, boxrule=1pt,
	left=10pt, right=10pt, top=10pt, bottom=10pt
}
\newtcolorbox{solutionbox}[1][]{
	colback=green!5!white, colframe=green!60!black,
	fonttitle=\bfseries, title=#1, arc=3pt, boxrule=1pt,
	left=10pt, right=10pt, top=10pt, bottom=10pt
}

% --- Macros ---
\newcommand{\F}{\mathbb{F}}
\newcommand{\C}{\mathcal{C}}
\newcommand{\Gal}{\text{Gal}}
\newcommand{\Aut}{\text{Aut}}

\begin{document}
	
	% ==========================================
	% PROBLEM
	% ==========================================
	\section*{Problem 1: Twisting a Code via Galois Descent}
	
	\begin{problembox}[The 1-Cocycle Condition over $\F_4 / \F_2$]
		Let $L = \F_4$ and $K = \F_2$. We construct $\F_4 = \F_2(\omega)$, where $\omega^2 + \omega + 1 = 0$. 
		The Galois group $\Gamma = \Gal(\F_4/\F_2)$ is cyclic of order 2, generated by the Frobenius automorphism $\sigma(x) = x^2$. Note that $\sigma(\omega) = \omega^2$ and $\sigma(\omega^2) = \omega$.
		
		Let $V = \F_4^2$ be the ambient space for our codes. A 1-cocycle $c: \Gamma \to \text{GL}_2(\F_4)$ is uniquely determined by the matrix it assigns to the generator $\sigma$. Let $M = c_\sigma$. 
		
		\vspace{0.5em}
		\textbf{Tasks:}
		\begin{enumerate}
			\item \textbf{The Cocycle Condition:} Prove that for $M$ to define a valid 1-cocycle, it must satisfy the matrix equation $M \sigma(M) = I_2$, where $I_2$ is the identity matrix and $\sigma(M)$ applies the Frobenius automorphism to each entry of $M$.
			\item \textbf{Verifying a Twist:} Let $M = \begin{bmatrix} 0 & 1 \\ 1 & 0 \end{bmatrix}$. Prove that $M$ defines a valid 1-cocycle.
			\item \textbf{Constructing the Descended Space:} By Galois descent, the ``twisted'' $\F_2$-vector space corresponding to this cocycle is defined as the fixed points of the twisted Frobenius action:
			$$ V_M = \{ v \in \F_4^2 \mid v M = \sigma(v) \} $$
			Find a specific basis for $V_M$ over $\F_2$. \textit{(Hint: Let $v = (a, b)$ where $a, b \in \F_4$. Set up the system of equations implied by $v M = (a^2, b^2)$ and solve for the allowable vectors).}
		\end{enumerate}
	\end{problembox}
	
	\vspace{1.5em}
	\noindent\textbf{Space for Calculations:}
	
	\newpage
	
	% ==========================================
	% SOLUTION KEY
	% ==========================================
	\section*{Solution Key: Twisting a Code via Galois Descent}
	
%	\begin{solutionbox}[Step-by-Step Derivation]
		\textbf{Part 1: The Cocycle Condition}\\
		By definition, a map $c: \Gamma \to \text{GL}_2(\F_4)$ is a 1-cocycle if it satisfies $c_{\tau \gamma} = c_\tau \tau(c_\gamma)$ for all $\tau, \gamma \in \Gamma$. \\
		Since $\Gamma = \langle \sigma \rangle$ has order 2, the identity element is $\sigma^2 = \text{id}$. \\
		Let $M = c_\sigma$. Evaluating the cocycle condition at $\tau = \sigma$ and $\gamma = \sigma$ gives:
		$$ c_{\sigma^2} = c_\sigma \sigma(c_\sigma) $$
		Since $\sigma^2$ is the identity, $c_{\sigma^2}$ must map to the identity matrix $I_2$. Substituting $M$ into the right side yields:
		$$ I_2 = M \sigma(M) $$
		Thus, the condition is necessary and sufficient for cyclic Galois groups of order 2.
		
		\vspace{1em}
		\textbf{Part 2: Verifying the Twist Matrix}\\
		We are given the permutation matrix $M = \begin{bmatrix} 0 & 1 \\ 1 & 0 \end{bmatrix}$. \\
		We apply the Frobenius automorphism $\sigma(x) = x^2$ to each entry. Since $0^2 = 0$ and $1^2 = 1$, the matrix is invariant under $\sigma$, meaning $\sigma(M) = M$. \\
		We check the cocycle condition:
		$$ M \sigma(M) = M M = \begin{bmatrix} 0 & 1 \\ 1 & 0 \end{bmatrix} \begin{bmatrix} 0 & 1 \\ 1 & 0 \end{bmatrix} = \begin{bmatrix} 1 & 0 \\ 0 & 1 \end{bmatrix} = I_2 $$
		The condition holds, so $M$ is a valid 1-cocycle.
		
		\vspace{1em}
		\textbf{Part 3: Constructing the Descended Space $V_M$}\\
		Let $v = (a, b)$ with $a, b \in \F_4$. The twisted Frobenius condition is $v M = \sigma(v)$. \\
		First, calculate $v M$:
		$$ (a, b) \begin{bmatrix} 0 & 1 \\ 1 & 0 \end{bmatrix} = (b, a) $$
		Next, calculate $\sigma(v)$:
		$$ \sigma(a, b) = (a^2, b^2) $$
		Equating the two yields the system:
		\begin{align*}
			b &= a^2 \\
			a &= b^2
		\end{align*}
		Substitute the first equation into the second: $a = (a^2)^2 = a^4$. \\
		In $\F_4$, the multiplicative group has order 3, so $x^3 = 1$ for all non-zero $x$, which means $x^4 = x$ holds for \textit{every} element in $\F_4$. Therefore, $a$ can be any element in $\F_4$. \\
		Since $b = a^2$, every vector in $V_M$ is of the form $(a, a^2)$ for $a \in \F_4$. 
		
		To find an $\F_2$-basis, we evaluate this form for basis elements of $\F_4$ over $\F_2$ (which are $1$ and $\omega$):
		\begin{itemize}
			\item Let $a = 1$: The vector is $(1, 1^2) = (1, 1)$.
			\item Let $a = \omega$: The vector is $(\omega, \omega^2)$.
		\end{itemize}
		These two vectors, $v_1 = (1, 1)$ and $v_2 = (\omega, \omega^2)$, are linearly independent over $\F_2$. \\
		Thus, the descended $\F_2$-vector space (the twisted code) has the basis:
		$$ \mathcal{B} = \{ (1, 1), (\omega, \omega^2) \} $$
		\textit{Note that while this space sits inside $\F_4^2$, it is strictly a 2-dimensional vector space over the base field $\F_2$, representing a unique isometric classification class descended from the extension field.}
%	\end{solutionbox}
	
\end{document}